\documentclass[a4paper,12pt,answers]{exam}
\usepackage{graphicx} % Required for inserting images
\graphicspath{ {images/} }
\usepackage{amsmath}
\usepackage{amssymb}
\usepackage{tabularx}
\usepackage{booktabs}   
\usepackage{xcolor} 
\usepackage{dsfont}
\usepackage{graphicx}  
\usepackage{caption} 
\usepackage{subcaption}

\title{CS215 Assignment-1}
\author{Malay Kedia, Aakash Gupta, Yash Singh}
\date{August 2024}

\begin{document}
\maketitle
\begin{questions}
\question{\bf Finding optimal bandwidth}\\
  
% \begin{center}
%     \includegraphics[width=0.75\linewidth]{3b.png}
%     \caption{\\Histogram plot}
% \end{center}
\textbf{Task 1}\\
\textbf{Part (a)}
\begin{solution}\\

We are asked to prove the following:

\[
\int \hat{f}(x)^2 dx = \frac{1}{n^2 h} \sum_{j=1}^{m} v_j^2
\]

where \( \hat{f}(x) \) is the histogram estimator, \( v_j \) is the number of points in the \( j \)-th bin, and \( h \) is the bin width.

\subsection*{Proof}

The histogram estimator \( \hat{f}(x) \) is defined as:
\[
\hat{f}(x) = \sum_{j=1}^{m} \frac{p_j}{h} \mathds{1}(x \in B_j)
\]
also, \[ p_j = \frac{v_j}{n}\]
so the expression becomes,
\[
\hat{f}(x) = \sum_{j=1}^{m} \frac{v_j}{nh} \mathds{1}(x \in B_j)
\]

where:
- \( v_j \) is the number of data points in the \( j \)-th bin.
- \( \mathds{1}(x \in B_j) \) is the indicator function that equals 1 if \( x \) belongs to bin \( B_j \) and 0 otherwise.


Next, we need to compute \( \hat{f}(x)^2 \):

\[
\hat{f}(x)^2 = \left( \sum_{j=1}^{m} \frac{v_j}{nh} \mathds{1}(x \in B_j) \right)^2
\]
Since the indicator functions \textbf{1}(x $\in B_j$)  are mutually exclusive (a point x can only belong to one bin at a time), we can ignore cross-terms in the expansion of the square, leading to,
\[
\hat{f}(x)^2 = \sum_{j=1}^{m} \left( \frac{v_j}{nh} \right)^2 \mathds{1}(x \in B_j)
\]

Now, we integrate \( \hat{f}(x)^2 \) over the entire domain of \( x \):

\[
\int \hat{f}(x)^2 dx = \int \sum_{j=1}^{m} \left( \frac{v_j}{nh} \right)^2 \mathds{1}(x \in B_j) dx
\]

Since the bins are disjoint, the integral over each bin \( B_j \) simplifies to the width of the bin, which is \( h \). Therefore, the integral becomes:

\[
\int \hat{f}(x)^2 dx = \sum_{j=1}^{m} \left( \frac{v_j}{nh} \right)^2 h
\]

Simplifying further:

\[
\int \hat{f}(x)^2 dx = \frac{1}{n^2 h} \sum_{j=1}^{m} v_j^2
\]

Thus, we have proven that:

\[
\int \hat{f}(x)^2 dx = \frac{1}{n^2 h} \sum_{j=1}^{m} v_j^2
\]

\end{solution}

\textbf{Part (b)}
\begin{solution}\\

We need to prove the following expression:

\[
\sum_{i=1}^n \hat{f}_{(-i)}(X_i) = \frac{1}{(n-1)h} \sum_{j=1}^m \left( v_j^2 - v_j \right)
\]

where \( \hat{f}_{(-i)} \) is the histogram estimator after removing the \( i \)-th observation, and \( v_j \) is the number of points in the \( j \)-th bin.

\subsection*{Proof}

The histogram estimator \( \hat{f}(x) \) for the complete dataset is given by:

\[
\hat{f}(x) = \sum_{j=1}^m \frac{v_j}{nh} \mathds{1}(x \in B_j)
\]

where \( v_j \) is the number of points in the \( j \)-th bin.

When we remove the \( i \)-th observation, the number of points in the bin containing \( X_i \) is reduced by 1. Therefore, the histogram estimator after removing the \( i \)-th point which maybe in any of the bin as indicated by introduction of predicate for the same, \( \hat{f}_{(-i)}(x) \), is:

\[
\hat{f}_{(-i)}(x) = \sum_{j=1}^m \frac{v_j - \mathds{1}(X_i \in B_j)}{(n-1)h} \mathds{1}(x \in B_j)
\]


Now, sum the values of \( \hat{f}_{(-i)}(X_i) \) for all \( i \):

\[
\sum_{i=1}^n \hat{f}_{(-i)}(X_i) = \sum_{i=1}^n \sum_{j=1}^m \frac{v_j - \mathds{1}(X_i \in B_j)}{(n-1)h} \mathds{1}(X_i \in B_j)
\]

For each bin \( j \), the indicator function \( \mathds{1}(X_i \in B_j) \) equals 1 if \( X_i \) is in bin \( j \), and 0 otherwise. Therefore, for each bin \( j \), the sum over \( i \) simplifies to:

\[
\sum_{i=1}^n \hat{f}_{(-i)}(X_i) = \frac{1}{(n-1)h} \sum_{j=1}^m \sum_{i \in B_j} (v_j - 1)
\]


For each bin \( j \), there are \( v_j \) points in total. Thus, the sum over \( i \in B_j \) yields \( v_j \) terms, each equal to \( v_j - 1 \). Hence, the inner sum simplifies to:

\[
\sum_{i \in B_j} (v_j - 1) = v_j(v_j - 1)
\]

Substituting this into the previous expression, we get:

\[
\sum_{i=1}^n \hat{f}_{(-i)}(X_i) = \frac{1}{(n-1)h} \sum_{j=1}^m v_j(v_j - 1)
\]

Expanding the product \( v_j(v_j - 1) \), we get:

\[
v_j(v_j - 1) = v_j^2 - v_j
\]

Thus, the final expression becomes:

\[
\sum_{i=1}^n \hat{f}_{(-i)}(X_i) = \frac{1}{(n-1)h} \sum_{j=1}^m \left( v_j^2 - v_j \right)
\]
\end{solution}
\textbf{Task 2}\\
\textbf{Part(a)}
\begin{solution}\\
    The probability distributions $p_j$ for all the bins as calculated in the code as $\frac{v_j}{n}$ are as follows:\\
    
% \end{solution}
% \begin{table}[h!]      
    % \caption{Probability Distribution Across Bins}  
    \begin{center}
        \begin{tabular}{|c|c|c|} 
        \hline
        \textbf{Bin Number} & \textbf{Bin Range} & \textbf{Probability} \\ 
        \hline
        Bin 1 & 0.00 - 0.40  & 0.0059 \\
        Bin 2 & 0.40 - 0.80  & 0.3824 \\
        Bin 3 & 0.80 - 1.20  & 0.3235 \\
        Bin 4 & 1.20 - 1.60  & 0.0529 \\
        Bin 5 & 1.60 - 2.00  & 0.1118 \\
        Bin 6 & 2.00 - 2.40  & 0.1000 \\
        Bin 7 & 2.40 - 2.80  & 0.0059 \\
        Bin 8 & 2.80 - 3.20  & 0.0000 \\
        Bin 9 & 3.20 - 3.60  & 0.0118 \\
        Bin 10 & 3.60 - 4.00 & 0.0059 \\
        \hline
    \end{tabular}
    \end{center}
% \end{table}
\end{solution}
\textbf{Part(b)}
\begin{solution}
    Bandwidth of the histogram is pretty large and number of bins is also less. Hence the histogram is \textbf{oversmooth}.
% \end{solution}
    
% \begin{figure}[htbp]
%     \centering

    % First image inside a minipage
    \begin{minipage}{0.50\textwidth}
        \centering
        \includegraphics[width=\linewidth]{10binhistogram.png}
        \caption{Oversmooth}
    \end{minipage}
    \hfill 
    % Second image inside a minipage
    \begin{minipage}{0.50\textwidth}
        \centering
        \includegraphics[width=\linewidth]{optimalhistogram.png} 
        \caption{optimal fit}
    \end{minipage}
% \end{figure}


\end{solution}
\textbf{Part(d)}
\begin{solution}
    Optimal bandwidth is obtained where cross validation score is minimum. From the curve, this is achieved when bandwidth is \textbf{around 0.08}.
    The exact optimal bandwidth from the code come out to be 0.083333 with 48 as number of optimal number of bins.
    
\begin{center}
\includegraphics[width=0.9\linewidth]{crossvalidation.png}
\caption{cross validation vs bandwidth curve}
\end{center}

\end{solution}
\end{questions}



% NEXT ANSWERS
\begin{questions}
\question{\bf Higher Order Regression}\\
\textbf{Part(a)}\\
\begin{solution}
    In a simple linear regression model, the equation of the least squares regression line is:
\[
\hat{y} = \hat{\beta}_0 + \hat{\beta}_1 x
\]
where symbols have their usual meanings.

The least squares estimates for $\hat{\beta}_1$ and $\hat{\beta}_0$ are given by:
\[
\hat{\beta}_1 = \frac{\sum_{i=1}^{n} x_i y_i-\bar{x}\sum_{i=1}^{n}y_i}{\sum_{i=1}^{n} x_i^{2} - n\sum_{i=1}^{n}\bar{x}^2}
\]
\[
\hat{\beta}_0 = \bar{y} - \hat{\beta}_1 \bar{x}
\]
where $\bar{x}$ and $\bar{y}$ are the means of $x$ and $y$, respectively.

We want to show that the point $(\bar{x}, \bar{y})$ lies on the regression line. This means that when $x = \bar{x}$, the predicted value $\hat{y}$ should be equal to $\bar{y}$.

Substitute $x = \bar{x}$ into the equation of the regression line:
\[
\hat{y} = \hat{\beta}_0 + \hat{\beta}_1 \bar{x}
\]

Now, substitute the expression for $\hat{\beta}_0$:
\[
\hat{y} = \left( \bar{y} - \hat{\beta}_1 \bar{x} \right) + \hat{\beta}_1 \bar{x}
\]

Simplifying this:
\[
\hat{y} = \bar{y} - \hat{\beta}_1 \bar{x} + \hat{\beta}_1 \bar{x}
\]
\[
\hat{y} = \bar{y}
\]

Thus, when $x = \bar{x}$, the predicted value $\hat{y}$ equals $\bar{y}$.

Hence, the point $(\bar{x}, \bar{y})$ always lies on the least squares regression line, as the predicted value $\hat{y}$ equals the mean value $\bar{y}$ when $x = \bar{x}$.
\end{solution}
\textbf{Part(b)}\\
\begin{solution}
    The least squares estimates for $\beta_0$ and $\beta_1$ can be derived by minimizing the sum of squared residuals:
\[
\text{Residual sum of squares (RSS)} = \sum_{i=1}^{n} (Y_i - \hat{Y_i})^2 = \sum_{i=1}^{n} (Y_i - (\hat{\beta}_0 + \hat{\beta}_1 x_i))^2
\]

The least squares estimates are:
\[
\hat{\beta}_1 = \frac{\sum_{i=1}^{n} x_i Y_i-\bar{x}\sum_{i=1}^{n}Y_i}{\sum_{i=1}^{n} x_i^{2} - n\sum_{i=1}^{n}\bar{x}^2}
\]
\[
\hat{\beta}_0 = \bar{Y} - \hat{\beta}_1 \bar{x}
\]
where $\bar{x}$ and $\bar{Y}$ are the means of $x$ and $Y$, respectively.

Now, suppose we use $z = x - \bar{x}$ as the regressor variable. The new model becomes:
\[
Y = \beta_0^* + \beta_1^* z + \epsilon
\]
Since $z = x - \bar{x}$, this new model can be rewritten as:
\[
Y = \beta_0^* + \beta_1^* (x - \bar{x}) + \epsilon
\]
Expanding this:
\[
Y = \beta_0^* + \beta_1^* x - \beta_1^* \bar{x} + \epsilon
\]
We can compare this to the original model $Y = \beta_0 + \beta_1 x + \epsilon$. Clearly, the coefficient $\beta_1^*$ is the same as $\beta_1$, i.e.,
\[
\beta_1^* = \beta_1
\]
Now, to match the original model, we can identify:
\[
\beta_0 = \beta_0^* - \beta_1 \bar{x}
\]
Thus, the intercept in the transformed model, $\beta_0^*$, relates to the intercept in the original model, $\beta_0$, as:
\[
\beta_0^* = \beta_0 + \beta_1 \bar{x}
\]

\textbf{Hence the following inferences can be drawn regarding relations between the least square estimates of the two models: }
\begin{itemize}
    \item The slope $\beta_1^*$ is the same as $\beta_1$: $\beta_1^* = \beta_1$.
    \item The intercept $\beta_0^*$ is shifted by $\beta_1 \bar{x}$: $\beta_0^* = \beta_0 + \beta_1 \bar{x}$.
\end{itemize}
\textbf{Differences between the two models are as follows:}
\begin{itemize}
    \item The new model, by centering the $x$ values (i.e., subtracting $\bar{x}$ from each $x_i$), makes the interpretation of the intercept $\beta_0^*$ simpler as $\beta_0^*$ represents the value of $Y$ when $x = \bar{x}$, i.e., when $z = 0$.
    \item In the original model, $\beta_0$ represents the expected value of $Y$ when $x = 0$, which might be less meaningful depending on the context (if 0 is outside the range of observed $x$ values).
\end{itemize}

Thus, the transformation to $z = x - \bar{x}$ simplifies the interpretation of the intercept but leaves the slope unchanged.
\end{solution}
\end{questions}
\end{document}
